\begin{table}[H]
\centering
\small
\renewcommand{\arraystretch}{1.6} 
\begin{tabularx}{\textwidth}{@{} l p{2.5cm} p{2.5cm} X X @{}} 
\toprule
\textbf{Framework} & \textbf{Core} & \textbf{Characteristic} & \textbf{Strength in Education} & \textbf{Limitation in Education} \\ \midrule

\textbf{LlamaIndex} & 
Data-Centric: Optimizing the interface between LLMs and private data. & 
Retrieval-heavy; focus on indexing and query engines.& 
Very powerful when it comes to multistep input output manipulation of tools, LLMs and states&
Often prioritizes data retrieval over complex, multi-step pedagogical reasoning. \\ \hline

\textbf{LangGraph} & 
Explicit state management via cyclic graphs. Supporting edges, nodes creation making full control of logic, branching& 
\textbf{Stateful, non-linear branching} with fine-grained control& 
Suitable for complex nonlinear task, with state\-aware and context, logic localization and independence &
Branching acts as a selling point and a potential logical error, when applying complex flow to simple tasks. Most of all, due to its high level abstraction, input output of LLM get unstable\\ \hline

\textbf{Google SDK} & 
Tool-Centric: Direct, low-latency access to model capabilities. & 
Functional/Linear; focus on tool-calling and result passing. & 
Deliver stable tool results, state and worker result if the pipeline is well designed&
Becomes unstable and difficult to maintain when including branching workflows. \\ \hline

\textbf{CrewAI} & 
Role-Centric: Simulating human-like collaborative teams. & 
Autonomous, task-based collaboration between predefined roles. & 
High level, well LLM and human interaction facilitation, capable of complex workflow&
Can lack the deterministic, full control orchestration required for strict academic knowledge modeling. \\ \bottomrule
\end{tabularx}
\caption{Comparison of LLM Orchestration Frameworks}\label{tab:frameworks}
\end{table}